% \begin{savequote}
% \begin{verbatim}
%  _________________________________________
% ( Your object is to save the world, while )
% ( still leading a pleasant life.          )
%  -----------------------------------------
%   o
%    o   \_\_    _/_/
%     o      \__/
%            (oo)\_______
%            (__)\       )\/\
%                ||----w |
%                ||     ||
% 
% \end{verbatim}
% \end{savequote}
\def\ChapterCode{TAE}
\chapter
[\CO{[TAE]} Tätigkeitsbild und Orientierung]
{Tätigkeitsbild und Orientierung}
\label{chp:TAE}

\ChapterIntro
{%
Ein Übersichtskapitel über das Umfeld und die Tätigkeiten in
den Sanitätsberufen.

\TODO{GABS}{2010-11-11}{Intro fehlt}{}
}
{%
\begin{description}
  \item[Maintainer:] Sebastian Gabriel
  \item[Autoren:] Diverse
  \item[Reviewer:] Standard-Reviewprozess
  \item[Version:] {\VarDocVersionTypeName} ({\VarDocVersionTypeDescription})
  \item[SHA1:] {\DocGitCommit}
\end{description}

{\TxTeilkapitelAASS}
}


\Dictum
[Albert Uderzo und René Goscinny: Asterix und Kleopatra]
{Ein großes Abenteuer in vielen bunten Bildern.}

\clearpage
\section{Präklinik}

\subsection{Sanitäter}

\subsubsection{Sanitäter in der Theorie}

Der Beruf bzw. die Tätigkeit des Sanitäters wurde 2002 durch
das neue \E{Sanitätergesetz} (SanG) neu geregelt. Demnach
unterteilt man in \E{Rettungssanitäter} oder
\E{Notfallsanitäter}. Der Beruf bzw. die Tätigkeiten des
Sanitäters dürfen nur nach Maßgabe des Sanitätergesetzes
 ausgeübt werden.

Sanitäter haben ihre Tätigkeit \E{ohne Ansehen der Person}
gewissenhaft auszuüben und haben das Wohl der Patienten und
der betreuten Personen nach
\E{Maßgabe der fachlichen und wissenschaftlichen Erkenntnisse und Erfahrungen}
zu wahren.
% Nötigenfalls ist ein Notarzt oder, wenn ein solcher nicht
% zur Verfügung steht, ein sonstiger zur selbständigen
% Berufsausübung berechtigter Arzt anzufordern.

\EE{Fortbildung} ist ein wichtiger Teil der Tätigkeit,
Sanitäter haben sich daher laufend tätigkeitsrelevant
fortzubilden.
Darüber hinaus müssen die \E{Berufs- bzw.
Tätigkeitspflichten} erfüllt werden
(z.\,B. Dokumentationspflicht, Verschwiegenheitspflicht,
Auskunftspflicht, vgl.
\FullPrettyRef{topic:sanitaeter-pflichten}).


% \clearpage % TEMP temp clearpage
\Layout{60}{Allgemeiner Tätigkeitsbereich}
{Der
Tätigkeitsbereich des Sanitäters beinhaltet eine Reihe von
Tätigkeiten, welche \EE{eigenverantwortlich} durchgeführt
werden. Dazu gehört die eigenverantwortliche Anwendung von
Maßnahmen der qualifizierten Ersten Hilfe, Sanitätshilfe und
Rettungstechnik, einschließlich diagnostischer und
therapeutischer Verrichtungen. } {}
{./images/TAE/dif-w2-zelt-1.jpg} {Fürsorge für erkrankte oder
hilfsbedürftige Personen: Die Kernaufgabe des
Sanitäters.} {}


\Layout{60}{Tätigkeitsbereich des Rettungssanitäters}{% Der
Der Tätigkeitsbereich des \EE{Rettungssanitäters} umfasst die
selbständige und \E{eigenverantwortliche Versorgung und
Betreuung kranker, verletzter und sonstiger hilfsbedürftiger
Personen}, die medizinisch indizierter Betreuung bedürfen,
vor und während des Transports, die Übernahme sowie die
Übergabe des Patienten oder der betreuten Person im
Zusammenhang mit einem Transport, eine qualifizierte
Durchführung von lebensrettenden Sofortmaßnahmen sowie die
sanitätsdienstliche Durchführung von Sondertransporten. } {}
{./images/TAE/tae-teamwork-transfer-lbk.jpg} {Die gute
Zusammenarbeit mit dem Spitalspersonal ist wichtig, wie zum
Beispiel bei dieser Überstellung eines Intensivpatienten:
Wenn danach alle Beteiligten danach so fröhlich sind ist das
ein gutes Zeichen.}
{}



\Text{Tätigkeitsbereich des Notfallsanitäters}
{Der
Tätigkeitsbereich des \EE{Notfallsanitäters} umfasst
zusätzlich die \E{Unterstützung des Arztes} bei allen
notfall- und katastrophenmedizinischen Maßnahmen
einschließlich der Betreuung und des sanitätsdienstlichen
Transports von \E{Notfallpatienten}, die Verabreichung von
dafür freigegebenen \E{Arzneimitteln}, die
eigenverantwortliche Betreuung der berufsspezifischen Geräte,
Materialien und Arzneimittel und die Mitarbeit in der
\E{Forschung}. Notfallsanitäter können die Berechtigung zur
Durchführung allgemeiner und besonderer
\E{Notfallkompetenzen} erwerben (Arzneimittellehre,
Venenzugang und Infusion, Beatmung und Intubation). }{}{}{}{}


\Layout{60}{Dienstverhältnis}
{Der Beruf bzw. die Tätigkeiten des Sanitäters
dürfen \E{ehrenamtlich}, \E{berufsmäßig}, als
\E{Zivildienstleistender} oder unter anderen bestimmten
Umständen\ZFootnote{Die Tätigkeiten des Sanitäters dürfen
ferner als Soldat im Bundesheer, als Organ des öffentlichen
Sicherheitsdienstes, Zollorgan, Strafvollzugsbediensteter,
Angehöriger eines sonstigen Wachkörpers} ausgeübt werden. Die
\EE{berufsmäßige} Ausübung von Tätigkeiten des Sanitäters
setzt neben der jeweiligen Fachausbildung die Absolvierung
eines \E{Berufsmoduls} voraus. 
}
{}
{./images/TAE/tae-okt26-alle.jpg} {Im Notfall ist die
Organisationszugehörigkeit nebensächlich.} {}


\Text{Befristung}{%
Die Berufs- und Tätigkeitsberechtigung ist mit jeweils zwei Jahren
\E{befristet}. Zur Verlängerung der Berufs- und Tätigkeitsberechtigung bedarf
es der Absolvierung von Fortbildungen sowie einer Rezertifizierung.
}{}{}{}{}


% (2) Notfallpatienten gemäß Abs. 1 Z 2 sind Patienten, bei denen im Rahmen einer
% akuten Erkrankung, einer Vergiftung oder eines Traumas eine lebensbedrohliche
% Störung einer vitalen Funktion eingetreten ist, einzutreten droht oder nicht
% sicher auszuschließen ist.

\subsubsection{Sanitäter in der Praxis}

Es ist derzeit noch relativ unklar wie der Sanitäter in das
etablierte Berufsfeld der Gesundheitsberufe \enquote{hineinpasst}.
Die Berufsgruppen der Ärzte, Hebammen, Pflegepersonal usw.
gibt es schon seit Jahrzehnten oder gar Jahrhunderten und es
blieb genug Zeit für Streitigkeiten bezüglich der jeweiligen
Kompetenz und der Abgrenzung der Tätigkeitsfelder. Des
öfteren hat der Gesetzgeber Regelungen aufgestellt und die
Gerichte Orientierungspunkte gegeben; darüber hinaus wurden
durch die ständige Arbeit miteinander Realitäten
\enquote{geschaffen}. In manchen Punkten ist aber auch heutzutage
einiges unklar, oder einem stetigen Wandel unterworfen.

\Text{Sanitätergesetz 2002}
{Mit der Schaffung des
Sanitätergesetzes im Jahr 2002 wollte man die Tätigkeit der
Sanitäter endlich klar regeln und Rechtssicherheit schaffen.
Gleich vorweg, wirkliche Klarheit über die Tätigkeiten des
Sanitäters wurde nicht geschaffen, wie unzählige Diskussionen
zu diversen Themen beweisen. Bezüglich der genauen
rechtlichen Regelungen sei auf das Kapitel \ref{chp:JUS}
verwiesen, an dieser Stelle wollen wir uns aus praktischer
Sicht dem Tätigkeitsfeld und den Aufgaben der Sanitäter
widmen. 
}{}{}{}{}

\Layout{60}{Das Miteinander}{%
Es scheint ein Naturgesetz zu sein dass, sobald zwei oder
mehr Berufsgruppen aufeinander treffen, es zu Streitigkeiten
bezüglich der jeweiligen Kompetenzen, Rechte und Befugnisse
kommt. Auf der Strecke bleibt oft das Potential welches
entwickelt werden könnte, würden die Berufsgruppen
konstruktiv und zielgerichtet zusammenarbeiten. In diesem
Sinne sind die \AASS{} auf ein interdisziplinäres
Miteinander der Berufsgruppen ausgerichtet.}
{}
{./images/TAE/tae-teamwork-arzt-sani.jpg}
{Miteinander statt gegeneinander: Arzt und Sanitäter bilden eine
Einheit.}
{GABS.}


\opt{STATUS-DRAFT}{
\Text{Der Sanitäter -- früher}{%
\TakeNotesLarge
}{}
{}
{}
{}
}

\Text{Der Sanitäter -- heute}{%
Das Aufgabengebiet des Sanitäter von heute wird immer mehr um
Tätigkeiten erweitert, die medizinisch-wissenschaftliche
Kenntnisse voraussetzen und noch vor einiger Zeit nur den
Ärzten oder anderen Berufsgruppen vorbehalten waren (z.\,B. 
Medikamentengabe).

\Z{Dies geschah anfangs durch das SanG, später aber auch
durch die Anforderungen, welche sich zwangsläufig aus der Arbeit in
der Praxis ergeben und z.\,T. gar nicht gesetzlich geregelt
sind.}

\Z{In letzter Zeit kommt es auch zu einem Einfluss privater,
meist kommerzieller, internationaler Ausbildungsprogramme wie
z.\,B. \Abk{AMLS}{\texttrademark}\ZFootnote{Advanced Medical
Life Support{\texttrademark}},
\Abk{ITLS}{\texttrademark}\ZFootnote{International Trauma
Life Support{\texttrademark}},
\Abk{PHTLS}{\texttrademark}\ZFootnote{Prehospital Trauma Life
Support{\texttrademark}} etc. Deren Lehrinhalte sind meist
auf den Kompetenzbereich von Paramedics\ZFootnote{Paramedics
sind gutausgebildete Sanitäter im englischsprachigen Raum,
welche in nicht-Notarzt-basierten Systemen eingesetzt
werden, d.\,h. sie haben in diesen Ländern das Personal mit
der höchsten medizinischen Ausbildung, welches im
präklinischen Regelbetrieb eingesetzt wird.} ausgelegt, und
reizen Grenzen dessen, was allgemein als Aufgaben der
Sanitäter angesehen wird, weitgehend aus, oder überschreiten
diese.} Diesbezüglich kommt es immer wieder zu Diskussionen,
was Sanitäter tun bzw.
unterlassen\footnote{Schädliche Dinge zu unterlassen ist
genauso wichtig wie nützliche Dinge zu tun.} \sout{dürfen}
\E{müssen}\footnote{Meist kann man es sich nicht aussuchen,
ob man eine Maßnahme ergreift oder nicht. Handlungen, die man
machen oder unterlassen darf und die das Patientenwohl
fördern, \E{müssen} i.\,d.\,R. entsprechend vorgenommen oder
unterlassen werden.}, und was sie \E{nicht} tun oder
unterlassen dürfen.
}{}{}{}{}

\Text{Kompetenzen werden erweitert}{%
Anfänglich war das Tätigkeitsprofil weitgehend auf das
traditionelle Gebiet der \enquote*{erweiterten Ersten Hilfe}
beschränkt. Die ändert sich zunehmend. Einerseits durch die
gesetzlich definierten Notfallkompetenzen, welche dem
ärztlichen Leiter in der großen Spielraum hinsichtlich der
Definition von Therapiemaßnahmen lassen, welche durch
(Notfall-) Sanitäter durchgeführt werden können;
andererseits werden andernorts auch Kompetenzerweiterungen
auch mit dem sich wandelnden Stand der Wissenschaft
begründet und durchgeführt. Als Beispiel sei die Einführung
des Larynxtubus genannt, für die es eine entsprechende
Rechtsansicht des zuständigen Bundesministerium gibt
\cite{BMG-92250_0080-II_A_2010_2010-12-15}. Andernorts wurde
bereits der Einsatz von Lachgas zur Schmerzbekämpfung
erprobt \cite{Hansak201207,Koppensteiner201208}.

Die Situation in Österreich ist jedoch noch recht
uneinheitlich: Während in manchen Regionen eifrig
Notfallsanitäter und Notfallkompetenzen ausgebildet bzw.
angewendet werden, sind manche (Landes-) Organisationen
äußerst zurückhaltend
\cite{Enne2012}.

}{}
{}
{}
{}


% \clearpage


\subsubsection{Anforderungen -- Beanspruchung -- Belastung: Was erwartet die
Welt von mir?}
\label{topic:taetigkeitsbild-sanitaeter-belastung}

\label{topic:belastung}
\label{topic:anforderung}

\Text{{\TxQuerverweise}}{%
\begin{inparaitem}
  \item Gesprächsführung: \CO{[AUP]} (\ref{chp:AUP});
  \item Belastungsreaktion, Burn-out:
  \prettyref{topic:ptss}, \prettyref{topic:burn-out}
\end{inparaitem}
}{}{}{}{}

% \begin{description}
%   \item[Anforderungen:] \Speech{fachliches Wissen ---
%   praktische Kompetenzen --- dieses in schwierigen Situationen umsetzen zu
%   können --- professionelle Haltung --- social/soft skills zwischenmenschliche
%   Fähigkeiten, Frustrationstoleranz --- Umgang mit mehrdeutigen
%   Situationen bzw. \enquote{schwierigem Klientel} --- Teamdifferenzen --- \ldots}
% \end{description}


% \begin{center}
% \ttfamily\small
% \begin{minipage}{45em}
% \begin{verbatim}
%  _________________________________________
% ( Your object is to save the world, while )
% ( still leading a pleasant life.          )
%  -----------------------------------------
%   o
%    o   \_\_    _/_/
%     o      \__/
%            (oo)\_______
%            (__)\       )\/\
%                ||----w |
%                ||     ||
% 
% \end{verbatim}
% \end{minipage}
% \end{center}

\TextSlide{Theoretische Anforderungen}{% 
Die Tätigkeit im medizinischen Bereich geht
aufgrund der Komplexität der Aufgaben mit einer Fülle von
körperlichen, wissensmäßigen und psychischen Anforderungen
einher. Einen nur kargen Hinweis darüber, welche Qualitäten
z.\,B. von einem Sanitäter erwartet werden, liefert das
Anforderungsprofil des Arbeitsmarktservice (AMS):
\begin{itemize}
\item Körperliche und psychische Belastbarkeit
\item rasches Auffassungs- und Reaktionsvermögen
\item Einfühlungsvermögen
\item Beobachtungsgabe
\item Kommunikations- und Teamfähigkeit
\item Zuverlässigkeit
\item Verantwortungsbewusstsein
\end{itemize}
}{%
\begin{itemize}
\item Körperliche und psychische Belastbarkeit
\item rasches Auffassungs- und Reaktionsvermögen
\item Einfühlungsvermögen
\item Beobachtungsgabe
\item Kommunikations- und Teamfähigkeit
\item Zuverlässigkeit
\item Verantwortungsbewusstsein
\end{itemize}
}{}{}{}



% Im Folgenden soll der Versuch unternommen werden diese
% nüchternen Auflistung von Begrifflichkeiten plastischer
% darzustellen und zu erweitern.

\TextSlide{Anforderungen und Belastungen in der Praxis}{%
Es versteht sich von selbst, dass die professionelle
Versorgung erkrankter bzw. verunfallter Menschen ein
medizinisches Wissen voraussetzt und damit auch einhergehend
die praktischen Kompetenzen, dieses \EE{zuverlässig und
eigenverantwortlich} in schwierigen Situationen umzusetzen zu
können.

Dringliche Versorgungssituationen bzw. Notfälle, in denen
jede zeitliche Verzögerung zulasten des Patienten fallen
könnten, machen ein \EE{schnelles Einschätzen der Lage} so
wie die Bereitschaft \EE{Verantwortung zu übernehmen}
notwendig. Da es sich stets um eine Interaktion zwischen zwei
Personen -- nämlich Versorgendem und zu Versorgenden --
handelt, gelten zwischenmenschliche Fähigkeiten wie höfliche
Umgangsformen, \EE{Kommunikationsfähigkeit} sowie
\EE{Einfühlungsvermögen} als Grundvoraussetzung.

Erkrankung bzw. Unfall stellen für die Betroffenen in
häufigen Fällen \E{kein alltägliches Geschehen} dar, sondern
haben oft einen kritisch-krisenhaften, jedenfalls aber
frustrierend-belastenden Charakter. Deshalb soll an dieser
Stelle die besondere Wichtigkeit einer \EE{einfühlenden und
verständnisvollen Grundhaltung} betont werden. Dies kann für
den professionellen Helfer auch damit einhergehen, \E{seine
eigenen Bedürfnisse} bzw. Empfindungen für den Moment
\E{hintanzustellen}, also eine hohe
\EE{Frustrationstoleranz}
erforderlich machen.

\Fazit{Der (richtige) Umgang mit Frustration ist wesentlich,
um im Rettungs- Krankentransportdienst auf Dauer \enquote{zu
überleben}.}

Im Rettungs- und Krankentransportdienst beschäftigte Personen
arbeiten
 die meiste Zeit in Teams, sodass Qualitäten wie Toleranz,
Kritikfähigkeit und Selbstdisziplin (\enquote{\EE{Teamfähigkeit}})
ebenso von hoher Bedeutung sind (Die Praxis zeigt, dass hier
häufig viel Spielraum für Verbesserungen besteht \Code{;-)}).

Weiters sei auf die Wichtigkeit der \EE{\E{körperlichen}
Belastungsfähigkeit} als erforderliches Kriterium
hingewiesen: das Tätigkeitsfeld des Rettungssanitäters
umfasst diverse Aufgaben, die mit körperlichem Einsatz und
dementsprechender Anstrengung verbunden sind.

Zu guter Letzt und im besonderen sei auf die Bedeutung
\EE{psychischer Belastbarkeit} hingewiesen, \AuthNotes{die
auch in den folgenden Kapiteln zum Thema gemacht werden
soll,} v.\,a. in Hinblick auf den Umgang mit psychischem
Stress:

Sanitätsdienst zu leisten, bedeutet tägliche
\E{Konfrontation} mit dem persönlichen Leid von Mitmenschen,
des weiteren die \E{Verantwortlichkeit} – ggfs. auch noch
unter Zeitdruck – darauf professionell zu reagieren und
zusätzlich die besonderen und generellen \E{Stressauslöser
des Arbeitsbereichs} eines Sanitäters zu meistern. Dies
benötigt eine hinreichend funktionierende psychische
Verarbeitungsfähigkeit, um ein gesundes seelisches
Gleichgewicht bzw. sein persönliches Wohlbefinden wahren und
pathologischen Entwicklungen (Burnout, Belastungsreaktion,
etc.) präventiv entgegenwirken zu können.

Weitere Informationen zu Stress und Belastungen finden sich
unter \FullPrettyRef{topic:belastung-psi}. 

}{%
\begin{itemize}
  \item Medizinisches Wissen und Kompetenz
  \item Zuverlässigkeit
  \item Eigenverantwortlichkeit
  \item Schnelles Einschätzen
  \item Verantwortung übernehmen
  \item Kommunikation
  \item Einfühlungsvermögen
  \item \EE{Frustrationstoleranz}
  \item Teamfähigkeit
  \item \E{körperliche} und \E{psychische} Belastungsfähigkeit
\end{itemize}
}{}{}{}



% \clearpage
\begin{WidePar}
\begin{multicols}{3}

\noindent\begin{minipage}{\linewidth}
\includegraphics[width=\linewidth]{./images/motal-aass/00800/traumacheck1.jpg}
\Captionof{figure}{\AuthNotesPicLegalOk{}{}Hilfeleistung.}
\end{minipage}\par

\noindent\begin{minipage}{\linewidth}
\includegraphics[width=\linewidth]{./images/bergung_bagger_00800.jpg}
\Captionof{figure}{\AuthNotesPicLegalOk{}{}Bergung unter
speziellen Bedingungen.}
\end{minipage}\par

\noindent\begin{minipage}{\linewidth}
\includegraphics[width=\linewidth]{./images/bergung_kran_00800.jpg}
\Captionof{figure}{\AuthNotesPicLegalOk{}{}Bergung unter
speziellen Bedingungen.}
\end{minipage}\par


\noindent\begin{minipage}{\linewidth}
\includegraphics[width=\linewidth]{./images/khd/gef-wind-rauch2}
\Captionof{figure}{Arbeiten im Gefahrenbereich.}
\end{minipage}\par

\noindent\begin{minipage}{\linewidth}
\includegraphics[width=\linewidth]{./images/motal-aass/intubation/dsc_4835-AASS-0030mm}
\Captionof{figure}{Assistieren bei (not-) ärztlichen Maßnahmen.}
\end{minipage}\par

\noindent\begin{minipage}{\linewidth}
\includegraphics[width=\linewidth]{./images/motal-aass/00800/neuro6}
\Captionof{figure}{Diagnostische Verrichtungen (Blutzuckermessung).}
\end{minipage}\par

\noindent\begin{minipage}{\linewidth}
\includegraphics[width=\linewidth]{./images/teaser/bh_ith_00800.jpg}
\Captionof{figure}{\AuthNotesPicLegalOk{}{}Ein
besonderer Arbeitsplatz: Intensivtransporthubschrauber.}
\end{minipage}\par

\noindent\begin{minipage}{\linewidth}
\includegraphics[width=\linewidth]{./images/TAE/dif-nacht-atmosphaere.jpg}
\Captionof{figure}{\AuthNotesPicLegalOk{}{}Hitze, Dunkelheit und
kontrolliertes Chaos: Ambulanzdienst auf dem Donauinselfest.}
\end{minipage}\par



\noindent\begin{minipage}{\linewidth}
\includegraphics[width=\linewidth]{./images/gabriel-sebastian-aass/DIF2010-000816-0154pt-0300dpi}
\Captionof{figure}{Gerüstet: Ambulanzraum bei einem Sanitätsdienst.}
\end{minipage}\par

\noindent\begin{minipage}{\linewidth}
\includegraphics[width=\linewidth]{./images/TAE/erste-hilfe-dif-schild-stange-2.jpg}
\Captionof{figure}{\AuthNotesPicLegalOk{}{}Licht, Menschenmassen und noch mehr
Hitze: Auch das ist das Donauinselfest.}
\end{minipage}\par


\noindent\begin{minipage}{\linewidth}
\includegraphics[width=\linewidth]{./images/gabriel-sebastian-aass/UE2011FORTUNA-00568-0154pt-0300dpi}
\Captionof{figure}{Führungsrolle bei der Bewältigung eines
Großschadensereignis: Der Leiter einer Sanitätshilfestelle
(SanHiSt).} \end{minipage}\par

\noindent\begin{minipage}{\linewidth}
\includegraphics[width=\linewidth]{./images/gabriel-sebastian-aass/UE2011FORTUNA-00518-0154pt-0300dpi}
\Captionof{figure}{Gefahrenstoffunfall.}
\end{minipage}\par

\noindent\begin{minipage}{\linewidth}
\includegraphics[width=\linewidth]{./images/gabriel-sebastian-aass/UEKAISERMUEHLENTUNNEL2011-000253-0154pt-0300dpi}
\Captionof{figure}{Zusammenarbeit zwischen den Organisationen.}
\end{minipage}\par

\noindent\begin{minipage}{\linewidth}
\includegraphics[width=\linewidth]{./images/gabriel-sebastian-aass/UE2011FORTUNA-00942-0154pt-0300dpi}
\Captionof{figure}{Organisation im Großschadensfall, hier
ein Wagenhalteplatz bei einer Nachtübung.}
\end{minipage}\par

\noindent\begin{minipage}{\linewidth}
\includegraphics[width=\linewidth]{./images/gabriel-sebastian-aass/UEKAISERMUEHLENTUNNEL2011-000193-0154pt-0300dpi}
\Captionof{figure}{Ungünstige Sichtverhältnisse, hier bei
einer Übung im Wiener Kaisermühlentunnel.} \end{minipage}\par

% \Layout{27}{}{
% 
% }
% {}
% {}
% {}
% {}
% \begin{figure}[H]
% \includegraphics[width=\linewidth]{./images/TAE/jacken-alt-grau.jpg}
% \Captionof{figure}{\AuthNotesPicLegalOk{}{}%
% % In den letzten zehn Jahren hat sich viel getan: 
% Die alten, grauen Jacken erscheinen heute wie ein Relikt aus grauer
% Vorzeit.}
% \end{figure}



% \Layout{27}{}{
% 
% }
% {}
% {}
% {}
% {}
\noindent\begin{minipage}{\linewidth}
\includegraphics[width=\linewidth]{./images/TAE/tae-okt26-bundesheer1.jpg}
\Captionof{figure}{\AuthNotesPicLegalOk{}{}Integraler Bestandteil der zivilen
und militärischen Landesverteidigung: Der Bundesheersanitäter.}
\end{minipage}\par


\noindent\begin{minipage}{\linewidth}
\includegraphics[width=\linewidth]{./images/gabriel-sebastian-aass/IMG_1853_1.jpg}
\Captionof{figure}{Die Zusammenarbeit mit der Exekutive -- auch an
sozialen Brennpunkten -- ist Alltag.}
\end{minipage}


% \Layout{27}{}{
% 
% }
% {}
% {}
% {}
% {}

\end{multicols}
\end{WidePar}


\section{Ärzte}


Der Arzt ist zur Ausübung der Medizin berufen. Dies umfaßt
jede auf medizinisch-wissenschaftlichen Erkenntnissen
begründete Tätigkeit, die unmittelbar am Menschen oder
mittelbar für den Menschen ausgeführt wird.\ZFootnote{Die
Ausübung der Medizin umfasst insbesonders die
Untersuchung auf das Vorliegen oder Nichtvorliegen von körperlichen und psychischen Krankheiten oder Störungen, die Beurteilung von diesen Zuständen bei Verwendung medizinisch-diagnostischer Hilfsmittel sowie deren Behandlung;
die Vornahme operativer Eingriffe einschließlich der Entnahme oder Infusion von Blut;
die Vorbeugung von Erkrankungen;
die Geburtshilfe;
die Verordnung von Heilmitteln, sowie
die Vornahme von Leichenöffnungen.
Ferner ist
der Arzt befugt, ärztliche Zeugnisse auszustellen und
ärztliche Gutachten zu erstatten.}





\Text
{Ausbildung}
{Die ärztliche Ausbildung gründet sich auf einem
abgeschlossenes Universitätsstudium, der weitere
Ausbildungsweg ist jedoch international sehr
unterschiedlich.}
{}
{}
{}
{}



\subsection{Situation in Österreich}

\Text
{Allgemeines}
{Der ärztliche Beruf wird durch das \I{Ärztegesetz} (\I{ÄrzteG}) geregelt.
Die selbstständige Ausübung des ärztlichen Berufes
ist ausschließlich \I[Arzt!für Allgemeinmedizin]{Ärzten für
Allgemeinmedizin}, \I[Arzt!approbierter]{approbierten
Ärzten} sowie \I[Facharzt]{Fachärzten}\index{Arzt!Fach-}
vorbehalten (\I{Arztvorbehalt}, von diesem Vorbehalt gibt es für andere medizinische Berufsgruppen
entsprechende Ausnahmeregelungen in den jeweiligen
Berufsgesetzen). \ZFootnote{Die in Ausbildung zum Arzt für
Allgemeinmedizin oder zum Facharzt befindlichen Ärzte (\I[Turnusarzt]{Turnusärzte}) sind lediglich zur
unselbstständigen Ausübung  in  als
Ausbildungsstätten anerkannten Einrichtungen, im Rahmen von
Lehrpraxen bzw. Lehrgruppenpraxen oder in Lehrambulatorien
unter Anleitung und Aufsicht der ausbildenden Ärzte
berechtigt.}

Interessensvertretung und Aufsichtsbehörde für alle Ärzte
ist die \I{Österreichische Ärztekammer}\index{Ärztekammer} bzw. die jeweils
zuständigen Landesärztekammern.}
{}
{}
{}
{}

\Text
{Ausbildung und Jus practicandi}
{In Österreich dauert das Studium der Humanmedizin je nach
Universität und Curriculum \Range{5}{6} Jahre. Im Anschluss an das
Studium verfügen die Absolventen \E{nicht} über das Recht zur
Berufsausübung, sie sind auch nach dem
Studium noch nicht als Ärzte, sondern als Mediziner zu
bezeichnen. 

Nach Abschluss des Studiums wird ein
Ausbildungsverhältnis angetreten, es erfolgt die Aufnahme in
die \EE{Ärzteliste}, welche von der Österreichischen
Ärztekammer geführt wird. Der Arzt in Ausbildung darf nun an
einer anerkannten Ausbildungseinrichtungen unter 
Aufsicht ärztlich tätig werden. Dies kann entweder im Rahmen
eines Ausbildungsverhältnisses zum Arzt für Allgemeinmedizin,
welches mindestens drei Jahre dauert, erfolgen oder in Ausbildung zu
einem Facharzt, diese Ausbildungszeit beträgt je nach Fach
\Range{5}{6} Jahre. Die Ausbildung zum Facharzt kann
grundsätzlich gleich nach dem Studium beginnen, bzw. auch erst nach der
Absolvierung der Ausbildung zum Arzt für Allgemeinmedizin
begonnen werden. Während der Ausbildungszeit durchläuft der
Auszubildende  im Turnus-System diverse
Fachabteilungen, welche in der jeweiligen Ausbildungsordnung
vorgegeben sind (daher die Bezeichnung
\T{Turnusarzt}\index{Arzt!Turnus-}). Die Turnus-Ausbildung
zum Arzt für Allgemeinmedizin dauert 3 Jahre bzw. zum
Facharzt je nach Fachgebiet \Range{5}{6} Jahre. 

Ein Arzt für Allgemeinmedizin verfügt über ein so genanntes
uneingeschränktes \T{Jus practicandi} (Recht zu
praktizieren). Ein Facharzt verfügt ebenfalls über ein Jus
practicandi, welches jedoch auf Tätigkeiten in seinem
Fachgebiet beschränkt ist, es sei denn, der betreffende
Facharzt hat auch die Ausbildung zum Arzt für
Allgemeinmedizin absolviert. Gegenwärtig gibt es in
Österreich über 40 Fachgebiete für die man sich
qualifizieren kann.}
{}
{}
{}
{}



\Text
{Der Arzt im Rettungs- und Krankentransportdienst}
{Das ÄrzteG sieht hierfür zwei Ausbildungen vor:
\begin{itemize}
  \item \T{Notarzt}: Die Ausbildung zum Notarzt setzt die
  Absolvierung eines 60-stündigen, von der österreichische Ärztekammer anerkannten und mit einer
  Prüfung abschließenden Diplomkurs
  voraus.
  Zum Erhalt der
  Tätigkeitsberechtigung ist eine regelmäßige
  Rezertifizierung (Notarzt-Refresher-Kurs) notwendig.
  \item \T{Leitender Notarzt}: Leitende Funktion, insbesonders 
  im Großschadensfall. Die Ausbildung zum Leitenden Notarzt setzt
  eine entsprechende Berufserfahrung als Notarzt sowie die
  Absolvierung eines entsprechenden Kurs voraus.
\end{itemize}

}
{}
{}
{}
{}







\section{Hebammen}

Der Hebammenberuf umfasst die Betreuung, Beratung und Pflege
der Schwangeren, Gebärenden und Wöchnerin, die
Beistandsleistung bei der Geburt sowie die Mitwirkung bei
der Mutterschafts- und Säuglingsfürsorge. Die Berufsgruppe
der Hebammen betreut werdende und gebärende
Mütter und ist zur Leitung von Geburten berechtigt.
Sie ist zur Durchführung von relevanten medizinischen
Maßnahmen, darunter fallen auch Blutabhnahmen, das Setzen von
peripheren Venenverweilkanülen, die Gabe von bestimmten
Medikamenten, sowie einschlägige chirurgische
Eingriffe (Versorgung eines Dammriss, \ldots) berechtigt.
Die Ausbildung zur Hebamme erfolgt in einer entsprechenden
Fachschule, neuerdings erfolgt die Ausbildung auch in
Fachhochschulen. Die Bezeichnung Hebamme gilt für weibliche
und männliche Berufsangehörige.
Die gesetzliche Grundlage bietet das \I{Hebammengesetz}
(\I{HebG}).

\section{Gehobener Gesundheits- und Krankenpflegedienst}

Zum gehobenen Gesundheits- und Krankenpflegedienst gehören
die diplomierten Gesundheits- und Krankenschwestern und
-pfleger (\I{DGKS},
\I{DGKP}),
die diplomierten Kinderkrankenschwestern und -pfleger 
(\I{DGKKS}, \I{DGKKP}), sowie die diplomierten
psychiatrischen Gesundheits- und
Krankenschwestern und -pfleger.
Die Ausbildung erfolgt
an Fachschulen (\enquote*{Schwesternschule}), aktuell
wird die Ausbildung auch im Rahmen eines Fachhochschulstudiums angeboten. 
Der Tätigkeitsbereich des gehobenen Gesundheits- und
Krankenpflegedienstes beinhaltet einen
eigenverantwortlichen 
Tätigkeitsbereich, welcher insbesonders die Pflege von
von Patienten vorsieht, sowie einen mitverantwortlich en
Tätigkeitsbereich, welcher die Durchführung von einfachen
ärztlichen Tätigkeiten (Blutentnahme, Verabreichung von
Infusionen und Medikamenten, \ldots) vorsieht. 
Für spezielle Tätigkeitsbereiche gibt es entsprechende
Sonderausbildungen (OP-Bereich, Anästhesie, Intensivpflege,
\ldots).


\ChapterEnd
